\chapter{Intraday Patterns}\label{ch:s1:intraday-patterns}

The US stock market trades from 9:30 to 4:00, but its futures trade around the clock, and Bondarenko and Muravyev found that four hours of the night,
around the European open, account for the entire average return of the E-mini S\&P 500; the other twenty hours average a noisy zero. The trading day has
its own shape too: volume piles up at the open and the close, the closing auction sets the prices that \omterm{def:s1:index-rebalancing:fund}{index funds} and dealers mark to, and in the last
half hour the market tends to continue the day's move, because dealers who are short gamma must hedge in its direction. This chapter decomposes the
day, simulates the hedging flow and the closing imbalance, and measures what each pattern is worth to trade. The build is \texttt{firm.intraday}.

\section{Overnight versus intraday returns}

\begin{definition}[Overnight return, intraday return]\label{def:s1:intraday-patterns:returns}
A security's \emph{overnight return}\index{overnight return} is its return from one day's close to the next day's open; its \emph{intraday
return}\index{intraday return} is the return from the open to the close. Their sum (in logs) is the close-to-close return.
\end{definition}

The split has been studied since the first intraday data: Cooper, Cliff and Gulen's ``like night and day'' paper is the classic reference, and Lou, Polk and
Skouras found that the profits of reversal and several momentum strategies were earned entirely overnight, with continuation within each part of the day
and reversal across them (chapter~2). How much of the market's own return is earned at night is harder to pin down than it looks. \texttt{firm.intraday}
simulates twenty years of an index in half-hour bars with a 7\% expected return spread in proportion to variance, a quarter of the variance arriving
overnight, and so a quarter of the expected return: 1.75\% a year overnight. Measured over the twenty years, the overnight share of the cumulative return is
54\% (4.4\% a year overnight, 3.7\% by day). The standard error of a twenty-year mean return with a volatility of 16\% is 3.6\% a year, and of its overnight
part 1.8\%: twenty years cannot tell a quarter from a half. Bondarenko and Muravyev's result is sharper because it rests on the futures' whole clock and
a Sharpe ratio of 1.6 in four hours, not on a split of a noisy mean.

\section{The volume smile and the auctions}

\begin{definition}[Volume smile]\label{def:s1:intraday-patterns:smile}
The \emph{volume smile}\index{volume smile} is the U-shaped pattern of trading volume (and volatility) across the trading day: high at the open, low at
midday, highest at the close, with the closing auction taking a large share.
\end{definition}

In the chapter's model the first and last half hours each carry 12.1\% of the day's volume and the midday half hour 4.9\%; a quarter of the volume
trades in the first and last half hours together. The close is where benchmark prices are set: \omterm{def:s1:index-rebalancing:fund}{index funds} (chapter~10), mutual funds that price their
shares at the close, and dealers who mark their books all trade there, through the closing auction (Book~1, chapter~13).

\begin{dated}{2026-09}{The Nasdaq closing cross}\label{dat:s1:intraday-patterns:nasdaq}
In Nasdaq's description of its closing cross (2019 edition of its FAQ), market-on-close, limit-on-close and imbalance-only orders are accepted before 3:50
p.m. Eastern Time; from 3:50 p.m. Nasdaq disseminates its Net Order Imbalance Indicator (every 10 seconds, then every second from 3:55 p.m.) and on-close
orders can no longer be cancelled or modified; market-on-close entry stops at 3:55 p.m., limit-on-close entry at 3:58 p.m., imbalance-only orders are
accepted until 4:00 p.m., when the closing process begins. Other exchanges run their own auctions with their own cut-offs.
\end{dated}

\begin{omfigure}
\begin{tikzpicture}
\begin{axis}[name=a,ybar,bar width=7pt,width=6.4cm,height=5.2cm,xlabel={\small half hour of the day},ylabel={\small share of volume (\%)},
  tick label style={font=\footnotesize},xtick={1,4,7,10,13},ymin=0,ymax=14,grid=major,grid style={gray!25},table/col sep=comma]
\addplot[fill=omDef!60,draw=omDef] table[x=bar,y=share_pct]{figdata/strategies-1/13-intraday-patterns/smile.csv};
\end{axis}
\begin{axis}[at={(a.east)},anchor=west,xshift=1.3cm,width=6.8cm,height=5.2cm,xlabel={\small return before the last half hour (\%)},
  ylabel={\small last half hour (bp)},tick label style={font=\footnotesize},grid=major,grid style={gray!25},legend pos=north west,
  legend style={font=\scriptsize},table/col sep=comma]
\addplot[black!50,thick,mark=o] table[x=rest0,y=last0_bp]{figdata/strategies-1/13-intraday-patterns/binned.csv};
\addplot[omThm,very thick,mark=*] table[x=rest5,y=last5_bp]{figdata/strategies-1/13-intraday-patterns/binned.csv};
\legend{no hedging flow,short-gamma hedging}
\end{axis}
\end{tikzpicture}

\omcaption{The chapter's day model. Left: share of daily volume by half hour. Right: mean return of the last half hour by decile of the return from the
previous close to the last half hour, without and with dealers hedging a short-gamma book. Data: \texttt{s1\_intraday}.}\label{fig:s1:intraday-patterns:smile}
\end{omfigure}

\section{End-of-day hedging flows}

\begin{definition}[End-of-day hedging flow]\label{def:s1:intraday-patterns:hedging}
An \emph{end-of-day hedging flow}\index{end-of-day hedging flow} is the trading that dealers and leveraged or inverse funds must do near the close to rebalance
their exposure after the day's move; when they are short gamma (Book~1, chapter~26), the hedge buys after rises and sells after falls, in the direction of
the move.
\end{definition}

Baltussen, Da, Lammers and Martens found market intraday momentum in over sixty futures on equities, bonds, commodities and currencies from 1974 to 2020: the
return of the last thirty minutes before the close is positively predicted by the return from the previous close to that point, significantly and in all
asset classes, and the effect reverts over the following days. They linked it to the gamma hedging of option market makers and leveraged ETFs. The
chapter's model plants that mechanism: dealers who are short gamma add 5\% of the day's move so far to the last half hour, and half of the push reverses
overnight.

\begin{center}
\small
\begin{tabular}{@{}lcc@{}}
\toprule
twenty synthetic years & no hedging flow & short-gamma hedging\\
\midrule
slope of the last half hour on the rest of the day; $t$ & $-0.006$; $-1.2$ & 0.044; 9.2\\
last-half-hour momentum: hit rate; Sharpe ratio after costs & 48.8\%; $-0.77$ & 54.1\%; 1.05\\
last-half-hour momentum: return a year & $-4.0\%$ & 5.5\%\\
closing-imbalance fade: Sharpe ratio; return a year & 1.09; 8.8\% & 1.08; 8.7\%\\
next \omterm{def:s1:intraday-patterns:returns}{overnight return} on the last half hour (slope) & $-0.048$ & $-0.053$\\
\bottomrule
\end{tabular}
\end{center}

Without hedging flow the last half hour is unpredictable (slope $-0.006$, $t = -1.2$), and a trade that follows the day's direction for the last half hour
loses its costs of one basis point a day: $-4.0\%$ a year. With it the slope is 0.044 ($t = 9.2$), the direction is right 54.1\% of the time, and the trade
earns 5.5\% a year at a Sharpe ratio of 1.05 (\cref{fig:s1:intraday-patterns:smile}, right). The trade is only as good as the hedging demand behind it:
when option dealers are long gamma (after investors sell options to them), the hedge trades against the move and the pattern reverses. Measuring the
dealers' net gamma is the practitioner's version of this chapter's \texttt{gamma} parameter.

\section{Close imbalances}

The closing auction publishes its imbalance before it prints. In the model an imbalance of one standard deviation pushes the close by ten basis points, and
60\% of the push reverses overnight. A trade that takes the other side of the published imbalance at the close and exits at the next open earns 8.8\% a year
at a Sharpe ratio of 1.09. The logic is chapter~10's in miniature: whoever must trade at the close pays whoever can wait until the open. The risk is also
the same: an imbalance may carry information, and the part that does not revert is the loss on the trades that should not have been made.

\section{Strategy files}

\begin{strategyfile}{Overnight drift}\label{strat:s1:intraday-patterns:overnight}
\sfield{Who pays you, and why} Investors who hold risk overnight, when uncertainty is resolved around the European open, and are paid for it.
\sfield{Instruments and venues} Index futures traded around the clock.
\sfield{Signal} The time of day: long in the hours around the European open.
\sfield{Sizing and execution} Enter before, exit after the window; small, repeated positions.
\sfield{Costs} Two futures trades a day.
\sfield{How it dies} As a risk premium it pays with losses in bad nights; crowding into the same window.
\sfield{Horizon, capacity, infrastructure} Hours; a round-the-clock futures desk.
\sfield{Backtest honestly} Futures prices at the exact times; costs; the COVID period in and out of the sample.
\sfield{Sources} Bondarenko and Muravyev (2022): the entire average return in four hours around the European open, Sharpe ratio 1.6.
\end{strategyfile}

\begin{strategyfile}{Last-half-hour momentum}\label{strat:s1:intraday-patterns:lasthalf}
\sfield{Who pays you, and why} Dealers and leveraged funds that must hedge in the direction of the day's move.
\sfield{Instruments and venues} Index and other liquid futures.
\sfield{Signal} The return from the previous close to 30 minutes before the close.
\sfield{Sizing and execution} In the direction of that return for the last half hour; flat at the close.
\sfield{Costs} One round trip a day in the most liquid contracts.
\sfield{How it dies} When dealers are long gamma; when others trade the same half hour first.
\sfield{Horizon, capacity, infrastructure} Half an hour; intraday data and execution.
\sfield{Backtest honestly} Prices at the exact bar boundaries; costs; the reversal the next day.
\sfield{Sources} Baltussen, Da, Lammers and Martens (2021); Gao, Han, Li and Zhou (2018, chapter~5).
\end{strategyfile}

\begin{strategyfile}{Close imbalance reversal}\label{strat:s1:intraday-patterns:imbalance}
\sfield{Who pays you, and why} Traders who must trade at the close and push its price.
\sfield{Instruments and venues} Stocks and their closing auctions.
\sfield{Signal} The published imbalance before the auction.
\sfield{Sizing and execution} Offset the imbalance in the auction; exit at the next open or through the next day.
\sfield{Costs} Auction fees; overnight risk.
\sfield{How it dies} Imbalances that carry information; competition to offset them.
\sfield{Horizon, capacity, infrastructure} Overnight; auction imbalance feeds and order-entry cut-offs.
\sfield{Backtest honestly} The imbalance as published at the cut-off, not the final one; fills in the auction.
\sfield{Sources} This chapter's simulation; the auction rules in the dated box.
\end{strategyfile}

\begin{strategyfile}{Hedging-flow timing}\label{strat:s1:intraday-patterns:hedging}
\sfield{Who pays you, and why} Dealers' mechanical hedging, when their gamma exposure is known or estimable.
\sfield{Instruments and venues} Index futures and options.
\sfield{Signal} An estimate of dealers' net gamma from option open interest, combined with the day's move.
\sfield{Sizing and execution} Scale the last-half-hour trade with the estimated gamma; reverse it when gamma is positive.
\sfield{Costs} As for last-half-hour momentum.
\sfield{How it dies} Wrong gamma estimates; changes in who is short gamma.
\sfield{Horizon, capacity, infrastructure} Intraday; option positioning data.
\sfield{Backtest honestly} Gamma estimates from data before each day.
\sfield{Sources} Baltussen, Da, Lammers and Martens (2021) on the link to gamma hedging.
\end{strategyfile}

\begin{strategyfile}{Open-auction reversal}\label{strat:s1:intraday-patterns:openreversal}
\sfield{Who pays you, and why} Overnight order flow that pushes the open, reversed during the day.
\sfield{Instruments and venues} Stocks and the opening auction.
\sfield{Signal} The \omterm{def:s1:intraday-patterns:returns}{overnight return}, against the stock's own history and its peers.
\sfield{Sizing and execution} Against the overnight move from the open, flat by the close.
\sfield{Costs} Wider spreads just after the open.
\sfield{How it dies} Continuation within the day; news behind the move.
\sfield{Horizon, capacity, infrastructure} Hours; opening auction data.
\sfield{Backtest honestly} Official opening prices; no signal from prices that could not be traded.
\sfield{Sources} Lou, Polk and Skouras (2019) on cross-period reversal; chapter~2.
\end{strategyfile}

\section{Tutorial: the market makes its money at night}\label{tut:s1:intraday-patterns:tut}

\begin{tutorial}
\textbf{Goal.} Simulate the day in half hours, with and without hedging flow, and measure the split of returns, the last-half-hour predictability and the
closing-imbalance trade. \textbf{End state:} the table and \cref{fig:s1:intraday-patterns:smile}.
\begin{tutsteps}
\item \textbf{The day}: overnight and half-hour returns, the hedging push and the imbalance.
\omcode{code/firm/intraday/firm_intraday.py}{44}{59}{The day model.}\label{lst:s1:intraday-patterns:day}
\item \textbf{The regression}: the last half hour on the rest of the day.
\omcode{code/firm/intraday/firm_intraday.py}{67}{74}{Last-half-hour predictability.}\label{lst:s1:intraday-patterns:reg}
\item \textbf{Run} \texttt{summary(0.0)}, \texttt{summary(0.05)} and \texttt{fig\_intraday.py}.
\end{tutsteps}
\textbf{What to change next.} Make dealers long gamma on some days and short on others, and give the trade a noisy estimate of which; run a hundred twenty-year
samples and plot the distribution of the measured overnight share.
\end{tutorial}

\section{Build: the trading day}\label{bld:s1:intraday-patterns:intraday}

\begin{build}
\bfield{Purpose} A half-hour model of the trading day with its \omterm{def:s1:intraday-patterns:smile}{volume smile}, hedging flow and closing imbalance, and the decompositions that study them.
\bfield{Interface} \texttt{DayConfig}, \texttt{simulate\_days(n, cfg, rng)}, \texttt{decompose(open\_, close)}, \texttt{last\_half\_hour(bars, overnight)},
\texttt{volume\_curve(bars, u)}.
\bfield{Rules} Expected returns in proportion to variance unless planted otherwise; the hedge and the imbalance act on the last half hour and revert
overnight.
\bfield{Acceptance tests} \texttt{code/firm/intraday/tests/}: the volume curve and the decomposition by hand; the variance split; no predictability without
hedging, and the planted slope with it.
\bfield{Stretch} Stock-level days with opening auctions; a gamma estimate from option open interest; Book~7, chapter~2's tape for the order-level version.
\end{build}

\begin{omsources}
\item O.~Bondarenko and D.~Muravyev, ``Market return around the clock: a puzzle'', \emph{Journal of Financial and Quantitative Analysis} 58(3), 2022 (online; issue of 2023).
\item D.~Lou, C.~Polk and S.~Skouras, ``A tug of war: overnight versus intraday expected returns'', \emph{Journal of Financial Economics} 134(1), 2019.
\item G.~Baltussen, Z.~Da, S.~Lammers and M.~Martens, ``Hedging demand and market intraday momentum'', \emph{Journal of Financial Economics} 142(1), 2021.
\item M.~J.~Cooper, M.~T.~Cliff and H.~Gulen, ``Return differences between trading and non-trading hours: like night and day'', working paper, 2008.
\item Nasdaq, Nasdaq Closing Cross FAQ.
\end{omsources}

\section{Exercises}

\begin{exercise}[$\star$]\label{exo:s1:intraday-patterns:1}
With an annual volatility of 16\%, what is the standard error of a twenty-year mean return? Of its overnight part, if a quarter of the variance is overnight?
\end{exercise}

\begin{exercise}[$\star$]\label{exo:s1:intraday-patterns:2}
If expected returns are proportional to variance and a quarter of the variance is overnight, how much of a 7\% annual return is earned overnight?
\end{exercise}

\begin{exercise}[$\star$]\label{exo:s1:intraday-patterns:3}
An imbalance of one standard deviation pushes the close ten basis points and 60\% reverts overnight. What does fading it earn, and what would earning that
every day give over a year?
\end{exercise}

\begin{exercise}[$\star\star$]\label{exo:s1:intraday-patterns:4}
Dealers are short gamma and the market is up 2\% by 3:30. With a hedging coefficient of 0.05, what push does the model add to the last half hour?
\end{exercise}

\begin{exercise}[$\star\star$]\label{exo:s1:intraday-patterns:5}
Why does a short-gamma hedge trade in the direction of the move, and a long-gamma hedge against it?
\end{exercise}

\begin{exercise}[$\star\star$]\label{exo:s1:intraday-patterns:6}
Why can twenty years of data not establish that most of the market's return is earned overnight, while Bondarenko and Muravyev's result is convincing?
\end{exercise}

\begin{exercise}[$\star\star\star$]\label{exo:s1:intraday-patterns:7}
\emph{Coding.} Run \texttt{summary(0.0)} and \texttt{summary(0.05)}; explain why the closing-imbalance fade is the same in both and the last-half-hour trade
is not.
\end{exercise}

\begin{exercise}[$\star\star\star$]\label{exo:s1:intraday-patterns:8}
\emph{Find the flaw.} ``\omterm{def:s1:intraday-patterns:returns}{Overnight returns} have been higher than daytime returns in our sample, so we hold the market only overnight and double our
Sharpe ratio.''
\end{exercise}

\section{Problem: The Market Makes Its Money at Night}

\begin{problem}[{Weekend problem --- the shape of the day}]\label{pb:s1:intraday-patterns:1}
The chapter's day model and the public record.

\medskip\noindent\textbf{Part I --- Night and day.}
\begin{enumerate}
\item Define overnight and \omterm{def:s1:intraday-patterns:returns}{intraday returns}.
\item What did Bondarenko and Muravyev, and Lou, Polk and Skouras, find?
\item What share of the model's return is planted overnight, and what share is measured?
\item Why do the two differ?
\end{enumerate}

\medskip\noindent\textbf{Part II --- The close.}
\begin{enumerate}[resume]
\item Define the \omterm{def:s1:intraday-patterns:smile}{volume smile} and give the model's shares.
\item Describe the closing cross's timeline.
\item Who trades at the close, and why?
\item What does the closing-imbalance fade earn?
\end{enumerate}

\medskip\noindent\textbf{Part III --- Hedging flow.}
\begin{enumerate}[resume]
\item Define \omterm{def:s1:intraday-patterns:hedging}{end-of-day hedging flow}.
\item What did Baltussen and co-authors find?
\item Give the last-half-hour slope and trade with and without hedging.
\item What happens to the pattern when dealers are long gamma?
\end{enumerate}

\medskip\noindent\textbf{Part IV --- The verdict.}
\begin{enumerate}[resume]
\item State the \emph{named result}: the share of the cumulative return earned overnight and the last-half-hour predictability with and without hedging
flow.
\item Which pattern is a risk premium and which is a liquidity trade?
\item How would you estimate dealers' gamma?
\item What costs matter most for these trades?
\item What would you check before trading the overnight window?
\item Which strategy file has the most capacity?
\item How do these patterns interact with chapter~10's index events?
\item In one sentence: what makes the close special?
\end{enumerate}
\end{problem}

\section{Interview questions}

\begin{interviewq}[$\star$ \iqroles{trader}]\label{iq:s1:intraday-patterns:1}
Why is volume highest at the open and the close?
\end{interviewq}

\begin{interviewq}[$\star\star$ \iqroles{researcher, trader}]\label{iq:s1:intraday-patterns:2}
What is market intraday momentum, and what might cause it?
\end{interviewq}

\begin{interviewq}[$\star\star$ \iqroles{trader}]\label{iq:s1:intraday-patterns:3}
How would you trade the closing auction's published imbalance?
\end{interviewq}

\begin{interviewq}[$\star\star$ \iqroles{researcher}]\label{iq:s1:intraday-patterns:4}
Most of the market's return seems to come overnight. How would you test whether that is real?
\end{interviewq}

\begin{interviewq}[$\star\star$ \iqroles{risk}]\label{iq:s1:intraday-patterns:5}
What risks does a strategy that holds positions only overnight carry?
\end{interviewq}

\begin{interviewq}[$\star\star\star$ \iqroles{researcher}]\label{iq:s1:intraday-patterns:6}
Dealers are short gamma $\Gamma$ in dollar terms and hedge at the close. Derive the size of their hedge trade after a market move $\Delta S$, and the price push
under a linear impact $\lambda$ per dollar traded.
\end{interviewq}
